\documentclass{RDJ}

\usepackage{fixltx2e}%For textsubscript
\usepackage{nicefrac}
\usepackage[symbol]{footmisc}
\usepackage{url}
\makeatletter
    \g@addto@macro{\UrlBreaks}{\UrlOrds}
\makeatother

\submissiondate{3 March 2017}%Editor will fill in this date
\reviseddate{7 July 2017}%Editor will fill in this date
\acceptdate{13 July 2017}%Editor will fill in this date
\pubYear{2017}%Editor will fill in this date
\pubVol{33}%Editor will fill in this numbers
\pubNum{1}%Editor will fill in this numbers
\setcounter{page}{42}%Editor will fill in this numbers

\title{Title 6-8 Words (Subsequent Single Line Running Header) [Book Antiqua 24 Bold Italic Centred - separator line 2 1/4 thick]}
\makeatletter
\let\thetile\@title
\makeatother
\RDJShortTitle{Title 6-8 Words}

\author{Author Times New Roman\textsuperscript{a}, 12 Bold\textsuperscript{b}, Italic Centred and Space Before 6\textsuperscript{c}}

\begin{document}

\maketitle

\RDJabstract{Abstract: Times New Roman 11 bold italic justified space before 12. Abstract times new Roman 11 bold italic justified space before 12. Abstract times new Roman 11 bold italic justified space before 12. Abstract times new Roman 11 bold italic justified space before 12. Abstract times new Roman 11 bold italic justified space before 12. Abstract times new Roman 11 bold italic justified space before 12. Abstract times new Roman 11 bold italic justified space before 12. Abstract times new Roman 11 bold italic justified space before 12. Abstract times new Roman 11 bold italic justified space before 12. Abstract times new Roman 11 bold italic justified space before 12. Abstract times new Roman 11 bold italic justified space before 12. Abstract times new Roman 11 bold italic justified space before 12.}

\begin{table}[b]
    \begin{tabularx}{\columnwidth}{lX}
        \hline
        \textsuperscript{a}& Membership Status. Affiliation. Address, address, address, Country. E-mail: times new Roman 10\\
        \textsuperscript{b}& Do not move this table. It is locked in this position and should resize as text is entered. Text around the table should update as required.\\
        \textsuperscript{c}& Membership Status. Affiliation. Address, address, address, Country. E-mail: times new Roman 10\\
    \end{tabularx}
\end{table}

\RDJkeywords{Times new Roman, 10, space before 12, not, in, title, or, abstract, no full stop at the end; Up to 5 keywords allowed.}

\RDJnomenclature{}

\RDJroman{
\begin{tabular}{ll}
a   &times new Roman 10 [unit]\\
A   &alphabetical order [unit]\\
B   &next symbol [unit]
\end{tabular}
}

\RDJgreek{
\begin{tabular}{ll}
$\alpha$    &times new Roman 10 [unit]\\
$\beta$     &alphabetical order [unit]\\
$\theta$    &next symbol [unit]
\end{tabular}
}

\RDJsubscripts{
\begin{tabular}{ll}
a&times new Roman 10 [unit]\\
A&alphabetical order [unit]\\
B&next symbol [unit]
\end{tabular}
}

\section{Arial 12 Bold Space Before 12 Heading 1}

Many formatting features of this journal have developed over decades of mostly the printed version and present policy is to honour this tradition by minimizing deviations from this formatting.   The word processing program of choice remains Microsoft Windows.   The editorial/copy editing process is integrated, so that as the final MS-Word version emerges, it merely needs to be converted to .pdf file format for publishing on the internet.

There is no indent after a heading.  All sections are numbered. Journal main text is in Times New Roman 10 font.  Figures, equation and tables are referred to using non-capitalised references and the numbers not in brackets.  Each figure and table must be at least so referred to before (or above) the figure (e.g. see figure \ref{fig:Figures}). The page format is double column, with a 1.67 cm all round margins, the columns are 8.5 cm wide and separated by 0.5 cm blank space.  Unfortunately this has implication for figures and tables, which must fit into a single column.   A figure or table which spans both columns must be the exception.

    \begin{figure}[htp!]
        \centering
        \includegraphics[width=7cm]{Fig1.pdf}
        \caption{Figure should fit into one column, titles should be brief, and may not exceed three indented lines in Arial 10, Space after 6, space before 6.}
        \label{fig:Figures}
    \end{figure}
    
Figures and photographs in colour are of course welcome.
Tables, such as table \ref{tab:Tables}, should have captions above the table. In general tables, figures and equation should be placed after the first reference to it in the text. All figures and tables must be referred to in the text. Leave an open line after a table.
    
\begin{table}[htp!]
	\caption{Table should fit into one column. Titles should be brief and may not exceed three indented lines in Arial 10, Space before 6, space before 6.}
	\label{tab:Tables}
	\centering
	\begin{tabular}{lll}
		\hline
		Column 1 & Column 2 & Column 3\\
		\hline
		text & text & text\\
		text & text & text\\
		\hline		
	\end{tabular}
\end{table} 

In general equations should be numbered. The numbers must be aligned to the right margin of the column, such as equation \ref{eq:Equation}. The equation is aligned to the left margin and not indented.

\begin{equation}
	(x+a)^n=\sum_{k=0}^nx^ka^{n-k} 
\end{equation}\label{eq:Equation}

Subsequent paragraphs are indented (hanging) in the first line by 0.5 cm, even if they start after a figure.  Referencing is carried out by means of a superscript number, author names may of course be stated as well, numbering and listing occurring in the order of citation.

\section{The Second Main Heading Arial 12 Bold}

\subsection{The first secondary heading Arial 11 bold, space before 12}

\subsubsection{The third level heading Arial 10 not bold, space before 12}

There may thus be up to three levels of embedded sections.

Referencing occurs via the number format in square brackets and the list of references following the main body of the paper.  A sample of the format of the references, more extensive than that provided in the instructions to authors is included below.   Do not use {\it et al}. in the list of references but always state all the authors. Always state page numbers or number of pages for papers.  The list of references appearing at the end of the paper is balanced over the two columns.  The authors name appear in italics, as do the Titles of the Journals.   Paper titles are not capitalised, book, theses and report titles are capitalized.

\section{Lists}

Here is an example bullet list:

\begin{itemize}
	\item Item 1
	\item Item 2
\end{itemize}

Here is an example numbered list:

\begin{enumerate}
	\item Item 1
	\item Item 2
\end{enumerate}

\section{References}

The list of references must be given at the end of the paper. References must be listed in the order that they appear in the main text. All references should be numbered. When citing a reference, put the reference number in square brackets, e.g. [1] or [1,4] if citing more than one reference..

Here are a couple of examples for references: Book \cite{04KR01,08BU01,12ME01}, PhD thesis \cite{00VA01}, Masters thesis \cite{00CO01}, journal \cite{09VA01}, website \cite{knuthwebsite}, conference \cite{97GO01}, British standard \cite{BS848}, in collection \cite{00KL01}.

Here is the Bibtex code for a British standard.

\begin{lstlisting}
@TechReport{BS848,
  title       = {Fans for General Purpo-
                ses. {P}erformance Test-
                ing Using Standardized 
                Airways},
  institution = {British Standards 
                Institution},
  year        = {1997},
  type        = {Standard},
  month       = {September},
  key         = {BS 848-1},
  timestamp   = {2017-12-11},
}
\end{lstlisting}

When using Bibtex, make sure there are spaces between the initials (in the case of more than one first name) to ensure all of them are displayed. Make sure there is a \lq\lq and\rq\rq  between each author name to have the correct punctuation in the reference list, otherwise there will not be a comma between the names of multiple authors. Full first names of authors are not allowed. Give all the initials.

\bibliography{sample.bib}

%\begin{thebibliography}{99}
%  \bibitem{AAR2004} {\em AAR}, \emph{Product Performance Profile for Pneumatic Dunnage}. 2004.
  
%  \bibitem{Venter2012} {\em Venter MP and Venter G}, Development and Validation of a Numerical Model for an Inflatable Paper Dunnage Bag, \emph{Packaging Technology and Science}, 25, 467--483, 2012.
  
%  \bibitem{Hallquist2006} {\em Hallquist JO}, LS-DYNA Theory Manual, Livermore Software Technology Corporation, Livermore, USA, 2006.
  
%  \bibitem{Suhey2005} {\em Suhey JD, Kim NH and Niezrecki C}, Numerical Modeling and Design of Inflatable Structures - Application to Open-Ocean-Aquaculture Cages, \emph{Aquaculture Engineering}, 2005, 33(4), 285--303.
  
%  \bibitem{Cerda2003} {\em Cerda E and Mahadevan L}, Geometry and Physics of Wrinkling, \emph{Physical Review Letters}, 90(7), 074302-1--074302-4, 2003.
  
%  \bibitem{Venter2016} {\em Venter MP and Venter G}, Simple Implementation of Plain Woven Polypropylene Fabric, \emph{Journal of Industrial Textiles}, 2016, Advance online publication. Doi:10.1177/1528083716665627.
  
%  \bibitem{Stander2012} {\em Stander N, Roux W, Goel T, Eggleston T and Craig K}, LS-OPT User's Manual - A Design Optimization and Probabilistic Analysis Tool for the Engineering Analyst, Livermore Software Technology Corporation, Livermore, USA, 2012.
  
%  \bibitem{Witowski2011} {\em Witowski K, Feucht M and Stander N}, An Effective Curve Matching Metric for Parameter Identification using Partial Mapping, 8th European LS-DYNA Users Conference, Strasbourg, 2011.
  
%  \bibitem{Venter2015} {\em Venter MP}, A Methodology for Numerical Prototyping of Inflatable Dunnage Bags, PhD Thesis, Faculty of Engineering at Stellenbosch University, Stellenbosch, South Africa, 2015.

%\end{thebibliography}


\end{document} 